\section{Primera iteración de la Fase 5. Análisis del Generador de cartuchos compactos}
\label{sec:F5Intro}

El PC ya ejecuta cartuchos completos: la autoría produce el GWI y el
simulador lo consume. Pero el GWI es un formato de \emph{autoría}:
JSON legible, imágenes PNG sueltas, ids en texto, un historial
cronológico de edición y campos pensados para que una persona (y el
diseñador gráfico) los lea y edite. El firmware del ESP32 no es una
persona ni un navegador: no tiene parser de JSON, no tiene sistema de
ficheros cómodo, y cada byte de RAM y cada ciclo cuentan.

Esta fase construye el \textbf{generador del cartucho compacto} (GWC):
la pieza que traduce el GWI, un fichero JSON editable con la fuente de datos, gráficos y sus relaciones, a un \textbf{fichero binario mínimo} que el firmware posterior puede lee en
una pasada y ejecuta sin interpretar nada. El GWC es al GWI lo que el
código objeto al código fuente: lo que queda cuando la validación, la
legibilidad y la reversibilidad ya no hacen falta, porque todo ese
trabajo lo ha hecho el PC.

El principio es el mismo que gobierna todo el proyecto: \emph{el cartucho define,
el motor ejecuta}; y su consecuencia práctica: \emph{el PC valida, el ESP32 confía}.

\section{Primera iteración: Análisis del generador del cartucho compacto}
\label{sec:F5I1}

Esta iteración analiza qué información debe viajar en el GWC, cómo se
transforma y con qué forma de fichero; las decisiones de realización
(código, validaciones, pruebas) son de la siguiente iteración.

\subsection{Por qué un formato binario}
\label{sec:F5Porque}

El firmware del ESP32 se escribirá en C++ sobre un microcontrolador
de recursos limitados. Un GWI del Parachute son decenas de KB de JSON
en texto y un ontón de PNG en un ZIP; leerlo exigiría al menos de un descompresor de ficheros ZIP, un parser JSON y un decodificador de ficheros PNG, todo en RAM y todo para acabar
teniendo los mismos datos que ya había. La decisión de diseño es
incluir en el GWC \textbf{los datos ya cocinados y en la forma en que el motor los consume}:

\begin{itemize}
	\item \textbf{Sin interpretación:} el firmware lee, no parsea. Los
	números están en binario con tamaños fijos; las referencias son
	índices de tabla, no cadenas \texttt{seg\_NNN} que haya que buscar.
	
	\item \textbf{Sin cálculo de escala:} decisión ya cerrada en la Fase
	2: las coordenadas y los sprites viajan ya escalados al tamaño real
	de la pantalla; el ESP32 los vuelca tal cual.
	
	\item \textbf{Sin validación:} todo el trabajo de verificación (los
	mismas comprobaciones bloqueantes del diseñador) se hace al generar;
	si el GWI no es coherente, el GWC no se produce.
	
	\item \textbf{Tamaño mínimo:} los sprites viajan como máscaras de
	1 bit (un LCD segmentado enciende píxeles, no colores), y el texto
	desaparece: los nombres legibles son fundamentales para las personas, pero el GWC
	ya no hay necesidad de editarlo.
\end{itemize}

\subsection{Inventario: qué viaja, qué se transforma, qué no viaja}
\label{sec:F5Inventario}

El punto de partida es el GWI real que produce la autoría. Sus
campos se clasifican en tres grupos:

% Definimos Y para que se repartan el espacio al 50% cada una (0.5 + 0.5 = 1)
\newcolumntype{Y}{>{\hsize=0.5\hsize}X}

\begin{small}
	\begin{xltabular}{\textwidth}{@{}l Y Y@{}}
		
	
		\toprule
		\textbf{Campo del GWI} & \textbf{Contenido} & \textbf{Destino en el GWC} \\
		\midrule
		\endfirsthead
		
		\multicolumn{3}{l}{\small\itshape (Continuación de la tabla anterior)} \\
		\toprule
		\textbf{Campo del GWI} & \textbf{Contenido} & \textbf{Destino en el GWC} \\
		\midrule
		\endhead
		
		\midrule
		\multicolumn{3}{r}{\small\itshape Continúa en la página siguiente...} \\
		\endfoot
		
		\bottomrule
		\endlastfoot
		
		\texttt{segmentos} & Tipos, posición, enlaces, propiedades de contadores y marca de primer jugador & \textbf{Viaja transformado}: tipos y propiedades como bytes; posiciones escaladas a la pantalla real; enlaces como índices binarios de segmento \\ \hline
		
		\texttt{sprites} & PNG en \texttt{segmentos/}, deduplicados & \textbf{Viaja transformado}: máscaras de 1 bit empaquetadas (8 píxeles por byte), escaladas al tamaño real \\ \hline
		
		\texttt{fin} & Mecanismo de fin (contador o familia de iconos) & \textbf{Viaja}: mecanismo y contador como bytes \\ \hline
		
		\texttt{secuencia} & Ticks de espera y lista de hasta cinco segmentos & \textbf{Viaja}: ticks y lista de índices \\ \hline
		
		\texttt{pausas} & \texttt{simulador} y \texttt{externo} & \textbf{Viaja}: la plataforma lee la suya; el simulador mantiene la suya para sus pruebas \\ \hline
		
		\texttt{aceleracion} & \texttt{cada}, \texttt{resta}, \texttt{minimo} & \textbf{Viaja} tal cual (enteros) \\ \hline
		
		\texttt{rutas} & \texttt{aumentarCada} & \textbf{Viaja} tal cual; el número de rutas lo dan los inicios del cartucho \\ \hline
		
		\texttt{estilo} & Estilo de máquina & \textbf{Viaja como comprobación}: el firmware comprueba que el cartucho es para su máquina (setup $\geq$ cartucho también en el hardware) \\ \hline
		
		\texttt{autoria} & Historial cronológico completo de edición & \textbf{No viaja}: es el mecanismo del editor; del GWC al GWI no hay camino de vuelta (el GWI es la fuente de verdad) \\ \hline
		
		\texttt{botones} & Botones usados del cartucho & \textbf{No viaja}: los puertos GPIO son de la máquina y viajan en el setup; el GWC es solo el juego \\ \hline
		
		\texttt{pantalla} & Referencias a \texttt{img\_f}, \texttt{img\_s}, \texttt{img\_m} & \texttt{img\_s} (origen de la extracción) y \texttt{img\_m} (portada) \textbf{no viajan}; \texttt{img\_f} queda como decisión abierta (sección de abajo) \\ \hline
		
		\texttt{nombre} & Título legible & \textbf{No viaja}: nadie lo muestra en juego; si el futuro menú de cartuchos lo pide, será una ampliación \\ \hline
		
		\texttt{Fantasmas} & Tipo 99 y segmentos con \texttt{visible: false} absorbidos. Restos informativos de las uniones & \textbf{No viajan}: nadie los enciende ni pinta (las barras y dígitos que el motor enciende sí viajan, con sus enlaces) \\ \hline
		
		\texttt{Enlaces redundantes} & Inversos calculados & \textbf{No viajan}: el motor no los lee; eran para la verificación y el panel \\ \hline
		
	\end{xltabular}

\nopagebreak
\captionof{table}{Campos GWI a GWC}
\label{tab:CamposGWIaGWC}
\end{small}

La lectura de la tabla es la definición misma del GWC: \emph{todo lo que el motor de estados 
y el motor gráfico leen en un tick, y nada más}.

\subsection{El escalado se hace al generar}
\label{sec:F5Escala}

Esto estuvo cerrado ya como criterio en la Fase 2: la resolución real de la pantalla final la declara el setup, y las coordenadas de los segmentos viven en el espacio de la imagen de
origen. En el simulador ambas cosas se reconcilian estirando al vuelo
con muestreo por vecino más cercano y escala por eje. Ese trabajo es
exactamente el que el ESP32 no debe hacer: el generador escala una
vez, al compilar el cartucho, y lo que viaja son las posiciones y
sprites definitivos en el espacio de la pantalla física. El firmware
recibe coordenadas y máscaras listas para volcar: cero multiplicaciones
en tiempo de juego.

Ya que un mismo GWI puede generar genera GWC distintos para
pantallas distintas, cuando cambia la resolución declarada en el setup. No se pudo hacer en el GWI para poder generar el cartucho para diversas resoluciones.

El GWI sigue siendo la fuente de verdad; el GWC es un \emph{build} para un objetivo concreto, como un binario compilado para una plataforma.

\subsection{Sprites compactos: máscaras de 1 bit}
\label{sec:F5Sprites}

El LCD de una Game \& Watch real es monocromo y segmentado: un
segmento está encendido o apagado. El GWC ya trae esa simplificación
de nacimiento — la deduplicación y las huellas píxel-exacta del
extractor son máscaras de bits —, así que el generador solo tiene que
empaquetarlas: cada sprite será un array de bytes con 8 píxeles por
byte, filas completas redondeadas a byte, precedido de su ancho y
alto. Un sprite de $16 \times 10$ ocupa 2 bytes por fila, 20 bytes en
total; los PNG de cientos de bytes han desaparecido.

Los sprites siguen deduplicados (la tabla de sprites es compartida
por los segmentos que la repiten, como en el GWI), y el orden es el
del GWI: determinista, mismo GWI $\rightarrow$ mismo GWC byte a
byte.

\subsection{La forma del fichero: tablas con conteos y registros fijos}
\label{sec:F5Forma}

El firmware necesita leer el GWC con el mínimo código posible. La
forma analizada sigue el patrón clásico de los formatos binarios
antiguos:

\begin{itemize}
	\item \textbf{Cabecera:} identificación (los cuatro bytes
	«GWC1»), versión del formato y suma de verificación\footnote{Suma de todos los valores módulo 65535 para que sea un dato entre 0000h y FFFFh, se cambia posteriormente a un CRC-16 que es mas adecuado, ver apartado \ref{sec:F5I14CRC}}. Es lo único
	que el firmware comprueba: la validación completa ya se hizo al
	generar.
	\item \textbf{Tablas con conteo y tamaño de registro fijo:} una
	tabla de sprites (direcciones de las máscaras) y una tabla de
	segmentos con todos los enlaces en cada registro; un enlace ausente
	es un valor reservado ($0\text{x}FFFF$). Registros de tamaño fijo
	significa acceso directo por índice: el motor dice «el segmento 23»
	y el firmware calcula su posición sin recorrer nada.
	\item \textbf{Bloques al final:} los parámetros del Director
	(pausas, aceleración, política de rutas), el fin de partida, la
	Secuencia y los bloques opcionales. Bloque marcado como ausente =
	comportamiento por defecto (pausa 500, sin acelerar, una ruta a la
	vez), igual que un GWI sin esos campos.
\end{itemize}

Con el Parachute (unos 40 segmentos, unos 45 sprites) la estimación
del resultado es de pocos KB: la cabecera y las tablas ocupan
apenas unos cientos de bytes, y las máscaras de todos los sprites
juntos caben con holgura en la RAM de un microcontrolador moderno.
El dato exacto será el resultado de la primera compilación real,
una cifra ilustrativa para la memoria de lo que cabe en el hardware
frente a lo que ocupaba el ZIP del GWI.

\subsection{El fondo: una decisión abierta que el formato no cierra}
\label{sec:F5Fondo}

El único bloque que en esta la Fase 5 no se decide es la imagen de
fondo (\texttt{img\_f}), porque depende de la presentación física
elegida en la Fase 6, deliberadamente aplazada a pruebas de
hardware:

\begin{itemize}
	\item Si la réplica física usa el diseño de la máquina original
	(arte impreso en una capa y LCD monocromo delante), el fondo no
	viaja: no hay nada que dibujar, el firmware solo enciende segmentos.
	\item Si usa un único TFT de color, el fondo debe estar disponible
	para el firmware, ya sea como bloque del GWC o como fichero
	acompañante en la memoria de programa.
\end{itemize}

El análisis deja la puerta abierta a ambas: el GWC reserva un
\textbf{bloque opcional de fondo} que la Fase 6 rellenará (o no)
según lo que midan las pruebas reales. El formato no se rompe por
ello: un bloque ausente es un bloque ausente.

\subsection{Los botones no viajan: el GWC es solo el juego}
\label{sec:F5Botones}

En el PC, el setup y el cartucho se cargan por separado y el
cartucho declara qué botones usa; en el hardware la relación es aún
más simple: los pines GPIO están cableados a la máquina y son
configuración del firmware (el setup del ESP32). El GWC no necesita
la botonera: el firmware lee sus pines, el motor traduce número
físico a movimiento, exactamente como hace hoy con la tabla del
simulador. La comprobación de compatibilidad (setup $\geq$
cartucho) se hace al generar: un GWC para una máquina que no tiene
sus botones no se produce; y por prudencia el GWC lleva el estilo
declarado, para que el firmware pueda rechazar un cartucho cambiado
a otra máquina con un error visible en lugar de un comportamiento
extraño.

\subsection{El generador confía en sí mismo; el firmware confía en el generador}
\label{sec:F5Confianza}

El reparto de responsabilidades es el que sigue:

\begin{itemize}
	\item \textbf{Al generar:} la verificación completa del diseñador,
	bloqueante: enlaces colgantes, contadores completos, jugador con
	movimiento, mecanismos de fin coherentes. Un GWI con errores no
	produce GWC, y el error es visible en la barra de estado, como
	siempre.
	
	\item \textbf{En el firmware:} solo la comprobación de identidad
	(magia, versión, suma). Todo lo demás es confianza: los índices son
	válidos, los enlaces apuntan a segmentos existentes, los contadores
	están declarados. Repetir la validación en el microcontrolador
	sería escribir el diseñador dos veces.
	
	\item \textbf{La reversibilidad no existe:} del GWC al GWI no hay
	camino de vuelta. La fuente de verdad sigue siendo el GWI editable;
	regenerar el GWC tras cualquier cambio es parte del flujo natural
	(compilar de nuevo), igual que el GWI se regenera desde el GWB.
	
\end{itemize}

\subsection{Estadística final}
\label{sec:F5Estadistica}
Al finalizar el proceso, la
generación presentará una comparación de tamaños con la ocupación
real de cada paso. Por un lado, la ocupación que tendrían todos los
sprites \emph{sin} deduplicación frente a la ocupación final de los
sprites en el GWC, con la cantidad ahorrada por la deduplicación.
Por otro, como comparación global, el tamaño del cartucho GWI
completo (el JSON con todos sus gráficos) frente al tamaño final
del cartucho GWC, con la cantidad total ahorrada. El Parachute real
será el caso de estudio: los números concretos ilustran en la
memoria la distancia entre un formato de autoría y uno de ejecución.


\subsection{Decisiones menores}
\label{sec:F5Menores}

\begin{itemize}
	\item La semilla y las constantes del PRNG no viajan: son
	constantes del firmware (las mismas 25173/13849 y la misma semilla
	del simulador), garantizando partidas idénticas entre PC y ESP32.
	
	\item La política de repeticiones de rutas
	(\texttt{MAX\_REPETICIONES\_RUTA}) sigue interna del motor: la
	variedad del generador no es dato del cartucho.
	
	\item El formato versiona desde el primer día (Tipo «GWC» Versión 1): las
	ampliaciones posibles, como el juego de tipo B con su cuarta ruta, la multi-colisión,
	más de una Secuencia, añadirán campos en futuras versiones, y el
	firmware comprueba la versión al cargarlo para verificar que sea un tipo correcto y soportado por el sistema.
	
	\item La generación se ofrecerá en el editor gráfico, ya que el editor tiene el cartucho en diseño en memoria, por lo que se añade un botón 
	«Generar GWI+GWC» junto al de verificar y guardar, que hará los dos procesos:
	primero guardar el GWI con su verificación previa necesaria para evitar errores, y si todo es
	correcto guardar el GWC, todo a la vez para que ambos cartuchos sean copia uno del otro siempre, y sin necesidad de salir del entorno.
	
	
\end{itemize}

Con este análisis cerrado, la iteración siguiente aborda la
realización: la estructura exacta de los registros en bytes, el
código del generador sobre el GWI en memoria y su verificación con
el cartucho real del Parachute.



